RF sensing can provide early evidence for cyberspace windows of superiority, but closed-set RF classifiers are unsafe when adversaries emit signal behaviors outside the trained taxonomy. We study open-world RF preliminary-action detection: recognizing known RF behaviors while routing novel behaviors as unknown. We introduce DQNGuard, a budgeted open-set decision layer over a multi-domain preliminary-action backbone. The backbone maps OTA RF windows into IQ, FFT, DCT, and polar representations; DQNGuard then uses predicted-class calibration, variance/energy-style guard evidence, and known-only thresholding to decide whether each window should be retained as known preliminary-action evidence or routed to an unknown-behavior pool. We evaluate over-the-air WiFi, Bluetooth, and Zigbee captures spanning five behaviors: Scan, Burst, Sustain, Hop, and Replay. Under fixed Scan-surrogate known-budget calibration, DQNGuard improves unknown-class detection relative to VarMax and a DQN-IDS-style confidence head while maintaining low known-sample rejection. A Target--Surrogate Matrix further shows that surrogate-open calibration is highly target-dependent: withheld classes are not interchangeable proxies for future unknown behaviors. These results position DQNGuard as an RF sensing and triage layer for QR-CWoS pipelines, where known preliminary actions can support downstream ATT\&CK/EW reasoning and rejected unknowns can feed label-making and continual-learning workflows.
